\section{Method}\label{sec:method}

\subsection{Problem setting}
An agent receives a prompt $x$ describing the task state and returns a proposal $a$. A validator $V(x,a)\in\{0,1\}$
returns whether $a$ is acceptable; in the CSTR it simulates the plant forward from the current state under the
proposed setpoints (Section~\ref{sec:setup}). In the baseline loop every proposal is validated: passing proposals are
executed, failing ones are returned with the validator's feedback for a new proposal, up to a budget. A
\emph{screen} maps signals $z(x,a)$ to an estimated failure risk $\hat r\in[0,1]$ and one of three decisions
(Fig.~\ref{fig:routing}):
\begin{equation}
d(\hat r)=\begin{cases}
\textsc{accept} & \hat r\le t_{\mathrm{acc}} \quad\text{(execute without validation)}\\
\textsc{reject} & \hat r\ge t_{\mathrm{rej}} \quad\text{(return to the agent without validation)}\\
\textsc{validate} & \text{otherwise.}
\end{cases}
\end{equation}
Both \textsc{accept} and \textsc{reject} skip a validator call. An unchecked accept of a failing proposal is the
safety-relevant error; an unchecked reject of a good proposal costs a regeneration.

\begin{figure}[t]
\centering
\begin{tikzpicture}[node distance=7mm and 9mm, every node/.style={font=\small},
  box/.style={draw, rounded corners, align=center, minimum height=8mm, inner sep=3pt},
  dec/.style={draw, diamond, aspect=2.2, align=center, inner sep=1pt}]
\node[box] (agent) {LLM agent\\proposes $a$};
\node[box, right=of agent] (sig) {signals $z(x,a)$:\\observables, token conf.,\\attention, hidden, grounding};
\node[box, right=of sig] (probe) {probe\\$\hat r$};
\node[dec, right=of probe] (route) {route};
\node[box, above right=3mm and 9mm of route] (exec) {execute};
\node[box, right=9mm of route] (val) {validator\\(digital twin)};
\node[box, below right=3mm and 9mm of route] (rej) {return to agent};
\draw[-{Latex}] (agent) -- (sig);
\draw[-{Latex}] (sig) -- (probe);
\draw[-{Latex}] (probe) -- (route);
\draw[-{Latex}] (route) |- node[pos=0.75, above]{\scriptsize $\hat r\le t_{\mathrm{acc}}$} (exec);
\draw[-{Latex}] (route) -- node[above]{\scriptsize else} (val);
\draw[-{Latex}] (route) |- node[pos=0.75, below]{\scriptsize $\hat r\ge t_{\mathrm{rej}}$} (rej);
\draw[-{Latex}] (val) -- node[right]{\scriptsize pass} (exec);
\draw[-{Latex}] (val) -- node[right]{\scriptsize fail + feedback} (rej);
\end{tikzpicture}
\caption{Skip-the-validator routing. A probe estimates the failure risk of each proposal from the chosen signals;
confident proposals are executed and clearly failing ones returned to the agent without a validator call. Under the
first-proposal-only rule (R0), retries are never accepted unchecked.}
\label{fig:routing}
\end{figure}

\subsection{Signals}
All internal signals are read from the forward pass that produces the proposal, at four relative depths of the
network (0.25, 0.5, 0.75 and 1.0).
\begin{itemize}
\item \textbf{Observables.} Task-state features visible in the prompt and features of the proposed action. In path
planning: graph size, query and path properties. In the CSTR: the measured temperature, level and actuator commands,
alarm flags, and the proposed setpoint changes. A \emph{plant-readings-only} variant omits the action features.
\item \textbf{Token confidence.} Log-probability, entropy and top-two margin statistics of the answer tokens, overall
and for numeric tokens.
\item \textbf{Attention by prompt region.} The share of attention mass from answer tokens to each prompt region
(instructions, knowledge graph or graph description, task state), with head-mean and head-max aggregation and
attention entropy.
\item \textbf{Hidden states.} The residual-stream vector at the last prompt token and its mean over answer tokens,
standardised and reduced to 64 principal components fitted on training data.
\item \textbf{Region-based attention grounding.} For each decision-relevant answer token, the share of attention on
the prompt region that decides it, relative to all task-content regions. In path planning, a step $u\to v$ is legal
only if $v$ appears in the adjacency line of $u$; the deciding region for the tokens of $v$ is that line. In the CSTR,
the deciding regions for each setpoint are the plant measurements that indicate its effect and the knowledge-graph
statements of the controller and balance equations that link it to the plant (pre-registered before computation).
Features are summarised over steps or setpoints (minimum, mean) and layers.
\end{itemize}
Combinations (\emph{all internals} = token confidence + attention + hidden states; \emph{observables + internals};
\emph{observables + grounding}) are fitted as single probes.

\subsection{Probes and thresholds}
Each signal set is fitted with an $\ell_2$-regularised logistic regression on the training partition (median
imputation and standardisation fitted on training data) and Platt-calibrated on a calibration partition. Thresholds
are chosen on the pooled development partitions $\mathcal D$ for a tolerance $X$ on the failure rate among unchecked
accepts. The accept threshold is the largest $t$ for which a one-sided Clopper--Pearson upper bound
\citep{clopper1934} at level $1-\delta_{\mathrm{sel}}$ ($\delta_{\mathrm{sel}}=0.05$) on the failure rate of
$\{a\in\mathcal D:\hat r(a)\le t\}$ is at most $X$. This bound-based rule replaced an initial point-estimate rule,
which proved too loose on certification data; the change was pre-registered before any affected test data were
analysed. The reject threshold is the smallest $t$ for which at least 95\% of $\{a\in\mathcal D:\hat r(a)\ge t\}$
fail. We report $X=10\%$ as primary and $X=5\%$ as secondary.

\subsection{Certification}
On an independent certification set used once, a screen is \emph{certified} at $X$ if the one-sided Clopper--Pearson
upper bound on the failure rate among its unchecked accepts, at level $1-\delta/K$ with $\delta=0.05$ and $K$ the
number of signal sets compared (Bonferroni), is at most $X$. A bound below $X$ requires a sufficiently large accept
set: with $K=11$ and no failures, at least about 52 unchecked accepts are needed at $X=10\%$, which was stated in
advance as a power limit for the 500-episode CSTR certification sets.

\subsection{Closed-loop evaluation and cost}
In closed loop each episode runs the propose--route--(validate)--retry loop with a budget of six proposals; if no
proposal is executed, the current setpoints are held. The executed outcome is the digital twin's verdict on the
executed setpoints with the plant's own future disturbances. We compare the screens with always-validate,
never-validate, and random routing at the combined probe's development accept and reject rates. Under the
\emph{first-proposal-only rule} (R0), a retry that the probe would accept is validated instead. The compute per
episode is
\begin{equation}
C = \sum_{\text{proposals}} \left(c_{\mathrm{gen}} + c_{\mathrm{feat}}\right) + n_{\mathrm{val}}\,c_{\mathrm{val}},
\end{equation}
where $c_{\mathrm{gen}}$ is generation time, $c_{\mathrm{feat}}$ the extra feature passes (zero for observables),
$n_{\mathrm{val}}$ the number of validator calls and $c_{\mathrm{val}}$ the validator cost per call. The screen's cost
does not depend on the simulator, whereas $c_{\mathrm{val}}$ grows with simulation horizon and fidelity. Because the
plant is paused during decisions, the logged decisions do not depend on $c_{\mathrm{val}}$, so $C$ can be recomputed
exactly for any validator cost, and the break-even cost $c^{*}_{\mathrm{val}}$ at which a screen matches
always-validate follows directly.
